%% kapa: exam=1 round=28 year=2017 period=1 form=A subject=economics subject_ko=경제학원론 seq=12 number=52
%% source: Q-Net 감정평가사 기출문제 (cst00309) artlSeq=5199603 file=extracted/1008/5199603/2017년 제28회 감정평가사 1차시험 문제지/1교시 A형.pdf
%% conversion: kordoc 4.24.51 -> tools/gen1.py (auto); review=proofread-sonnet; flags=none
\begin{question}{52}
다음 ( )안에 들어갈 내용으로 알맞은 것은?
\begin{framedtext}
관상용 나무재배는 공기를 정화하는 긍정적 외부효과(externality)를 발생시킨다. 나무재배 시 사회적 효용은 사적 효용보다(과) ( ㄱ ), 사회적 최적 재배량은 사적 균형 재배량보다(과) ( ㄴ ).
\end{framedtext}
\begin{choices}
  \item \makebox[7em][l]{ㄱ: 크며}ㄴ: 많다
  \item \makebox[7em][l]{ㄱ: 크며}ㄴ: 적다
  \item \makebox[7em][l]{ㄱ: 작으며}ㄴ: 많다
  \item \makebox[7em][l]{ㄱ: 작으며}ㄴ: 적다
  \item \makebox[7em][l]{ㄱ: 동일하고}ㄴ: 동일하다
\end{choices}
\end{question}
\answer{1}
